\section{Step 12 Keep the complete product for checking}\label{expression-s12}

Use \texttt{dmul} rather than the word-result multiply. Two signed-word operands have a product that fits in a signed doubleword, including the minimum value multiplied by itself. A word multiply would discard the high product bits before the evaluator could decide whether the result overflowed.

\Needspace{10\baselineskip}
\subsection{Write or inspect}
Replace the temporary return in the \texttt{apply\_multiply} region. Preserve its surrounding labels and frame scaffolding.

\Needspace{4\baselineskip}
\MARSCodeLabel{MIPS assembly}
\begin{minted}[samepage=true]{mars-idea}
    dmul $t0, $a0, $a1
\end{minted}
\subsection{Expected observation}
46341 times 46341 leaves exact candidate 2147488281, which exceeds the signed-word maximum.

\Needspace{11\baselineskip}
\subsection{Check your understanding}
Why is a word multiply insufficient here?

\begin{itemize}
\item \textbf{A} A word multiply always traps when its signed result is too large.
\item \textbf{B} It treats every operand as unsigned.
\item \textbf{C} It can discard the evidence needed to detect signed-word overflow.
\end{itemize}
